% ---------------------------------------------------------------------
%  CHƯƠNG 2 — QUY CÁCH TRÌNH BÀY
%  Chương này đóng vai trò sổ tra cứu: mỗi mục là một mẫu em copy rồi sửa.
%  Khi viết đồ án thật, em thay cả chương này bằng nội dung của mình.
% ---------------------------------------------------------------------
\chapter{Quy cách trình bày trong báo cáo}
\label{ch:soanthao}

\begin{muctieu}
Chương này gom sẵn mẫu cho mọi thành phần em cần dùng: hình, bảng, công
thức, mã nguồn, giải thuật, trích dẫn, chú thích và các hộp nhấn mạnh. Mỗi
mục có đoạn mã mẫu, em copy rồi thay nội dung.
\end{muctieu}

\begin{luuy}
Khi bắt đầu viết đồ án thật, hãy xoá cả chương này. Nó chỉ là sổ tra cứu,
không phải một chương của báo cáo.
\end{luuy}

\section{Chèn hình ảnh}

Đặt file ảnh vào thư mục \code{figs/} rồi dùng mẫu dưới đây. Mọi hình đều
phải có chú thích, có nhãn, và phải được nhắc tới trong thân bài như
Hình~\ref{fig:vidu}.

\begin{figure}[H]
\centering
\includegraphics[width=0.72\textwidth]{dau-campus.png}
\caption{Một ví dụ về chèn hình ảnh từ file.}
\label{fig:vidu}
\end{figure}

Nếu hình là sơ đồ do em tự vẽ thì nên vẽ thẳng bằng TikZ như
Hình~\ref{fig:luongxuly}, vì sơ đồ vẽ bằng TikZ luôn nét khi phóng to và
dùng đúng bộ màu của trường.

\begin{figure}[H]
\centering
% \resizebox ép sơ đồ vừa đúng bề ngang trang, tránh tràn lề
\resizebox{\textwidth}{!}{%
\begin{tikzpicture}[node distance=8mm and 12mm]
  \node[io] (in) {Ảnh đầu vào};
  \node[process, right=of in] (pre) {Tiền xử lý};
  \node[decision, right=of pre] (chk) {Đủ chất lượng?};
  \node[process, right=of chk] (mod) {Mô hình CNN};
  \node[startstop, below=of chk] (loai) {Loại bỏ};
  \draw[flowarr] (in)  -- (pre);
  \draw[flowarr] (pre) -- (chk);
  \draw[flowarr] (chk) -- node[above, font=\scriptsize] {có} (mod);
  \draw[flowarr] (chk) -- node[right, font=\scriptsize] {không} (loai);
\end{tikzpicture}}
\caption{Một ví dụ sơ đồ luồng xử lý vẽ bằng TikZ.}
\label{fig:luongxuly}
\end{figure}

\section{Chèn bảng}

Bảng dùng chú thích đặt phía trên, đúng quy cách trình bày khoa học. Dùng
\code{booktabs} với ba loại đường kẻ \code{toprule}, \code{midrule},
\code{bottomrule}, tránh kẻ dọc cho mọi ô.

\begin{table}[H]
\centering
\caption{Một ví dụ về bảng so sánh kết quả.}
\label{tab:sosanh}
\begin{tabular*}{\textwidth}{@{\extracolsep{\fill}}>{\bfseries\color{daunavy}}l cccc}
\toprule
\rowcolor{gray100}\bfseries\color{daunavy}Mô hình & \bfseries\color{daumaroon}Accuracy & \bfseries\color{daumaroon}Precision & \bfseries\color{daumaroon}Recall & \bfseries\color{daumaroon}F1 \\
\midrule
Huấn luyện từ đầu   & 0,848 & 0,841 & 0,836 & 0,838 \\
Học chuyển giao     & 0,912 & 0,907 & 0,903 & 0,905 \\
\bottomrule
\end{tabular*}
\end{table}

\begin{luuy}
Bảng rộng quá lề là lỗi hay gặp nhất. Nếu bảng tràn, hãy giảm cỡ chữ bằng
\code{small} hoặc \code{footnotesize}, hoặc gộp bớt cột. Đừng để LaTeX tự
tràn rồi bỏ qua.
\end{luuy}

\section{Chèn công thức toán học}

Công thức đặt giữa dòng và có đánh số thì dùng môi trường \code{equation},
như Công thức~\eqref{eq:softmax}. Công thức ngắn chèn trong câu thì viết
giữa hai dấu đô la, ví dụ $y = f(x)$.

\begin{equation}
p(y = k \mid \mathbf{x}) = \frac{\exp(z_k)}{\sum_{j=1}^{K} \exp(z_j)},
\label{eq:softmax}
\end{equation}

trong đó $z_k$ là điểm số thô mà mạng gán cho lớp $k$, và $K$ là số lớp. Mọi
ký hiệu xuất hiện trong công thức đều phải được giải thích ngay sau đó, như
câu này đang làm với \eqref{eq:softmax}.

\begin{ghinho}
Công thức là một phần của câu văn, không phải một khối trôi nổi. Trước công
thức phải có câu dẫn vào, sau công thức phải có câu giải thích ký hiệu hoặc
ý nghĩa. Công thức đứng một mình không ai đọc.
\end{ghinho}

\section{Chèn mã nguồn}

Mẫu này dùng gói \code{listings} thay cho \code{minted} của bản gốc, vì
\code{listings} không cần bật \code{-shell-escape} nên build được ngay trên
Overleaf lẫn máy cá nhân.

Cách thứ nhất là gõ trực tiếp, kết quả như Mã nguồn~\ref{lst:tructiep}.

\begin{lstlisting}[language=Python, caption={Chèn mã nguồn trực tiếp vào báo cáo.}, label={lst:tructiep}]
import torch.nn as nn

# A small convolutional block: conv, batch norm, activation
def conv_block(in_ch, out_ch):
    return nn.Sequential(
        nn.Conv2d(in_ch, out_ch, kernel_size=3, padding=1),
        nn.BatchNorm2d(out_ch),
        nn.ReLU(inplace=True),
    )
\end{lstlisting}

Cách thứ hai là nạp mã từ một file có sẵn trong đồ án. Cách này tiện vì báo
cáo luôn khớp với mã thật, sửa mã là báo cáo tự cập nhật theo. Mã
nguồn~\ref{lst:tufile} được nạp từ file \code{code/vi-du-huan-luyen.py}.

\lstinputlisting[language=Python, caption={Chèn mã nguồn từ file có sẵn.}, label={lst:tufile}]{code/vi-du-huan-luyen.py}

\begin{luuy}
Quy ước của khoa: comment trong code viết bằng tiếng Anh, còn văn xuôi của
báo cáo viết bằng tiếng Việt. Không đặt tiếng Việt bên trong khối mã nguồn,
vì phông mã nguồn không có đủ dấu tiếng Việt.
\end{luuy}

Chỉ đưa vào thân bài những đoạn mã cốt lõi. Mã dài để ở phụ lục hoặc chỉ dẫn
tới kho \code{git} của đồ án.

\section{Biểu diễn giải thuật}

Giải thuật trình bày bằng mã giả, không phải mã của một ngôn ngữ cụ thể,
để người đọc nắm được ý tưởng mà không cần biết ngôn ngữ đó.
Giải thuật~\ref{alg:nguyento} là một ví dụ.

\begin{algorithm}[H]
\caption{Kiểm tra một số tự nhiên có phải số nguyên tố hay không}
\label{alg:nguyento}
\hspace*{\algorithmicindent}\textbf{Đầu vào:} số tự nhiên $n$\\
\hspace*{\algorithmicindent}\textbf{Đầu ra:} \textsc{True} nếu $n$ là số nguyên tố, ngược lại \textsc{False}
\begin{algorithmic}[1]
\Function{LaSoNguyenTo}{$n$}
  \If{$n < 2$}
    \State \Return \textsc{False}
  \EndIf
  \For{$i \gets 2$ \textbf{to} $\lfloor \sqrt{n} \rfloor$}
    \If{$n \bmod i = 0$}
      \State \Return \textsc{False}
    \EndIf
  \EndFor
  \State \Return \textsc{True}
\EndFunction
\end{algorithmic}
\end{algorithm}

\section{Trích dẫn tài liệu tham khảo}

Tài liệu tham khảo khai báo trong file \code{refs.bib}, mỗi tài liệu một
khối. Trong thân bài em gọi bằng lệnh \code{cite} kèm khoá của tài liệu, ví
dụ mô hình AlexNet~\cite{Krizhevsky2012} và mạng ResNet~\cite{He2016}. LaTeX
tự đánh số và tự dựng danh mục ở cuối báo cáo, nên em không phải đánh số tay.

\begin{ghinho}
Chỉ liệt kê trong \code{refs.bib} những tài liệu em thực sự trích dẫn. Tài
liệu khai báo mà không gọi \code{cite} sẽ không xuất hiện trong danh mục.
\end{ghinho}

\section{Chú thích cuối trang}

Khi cần giải thích thêm mà không muốn làm gãy mạch câu, dùng lệnh
\code{footnote}\footnote{Đây là một dòng chú thích cuối trang.}. Đừng lạm
dụng, mỗi trang nhiều lắm một đến hai chú thích.

\section{Các hộp nhấn mạnh}

Mẫu này có sẵn bốn loại hộp để làm nổi những ý quan trọng.

\begin{dinhnghia}[Học chuyển giao]
Học chuyển giao \en{(transfer learning)} là kỹ thuật tận dụng tri thức mà mô
hình đã học được trên một bài toán nguồn có nhiều dữ liệu, để giải một bài
toán đích có ít dữ liệu hơn.
\end{dinhnghia}

\begin{ghinho}
Hộp ghi nhớ dùng cho ý chốt mà em muốn người đọc mang theo sau khi đọc xong
mục đó.
\end{ghinho}

\begin{luuy}
Hộp lưu ý dùng cho cảnh báo, điều kiện áp dụng, hoặc lỗi mà người đọc dễ mắc
nếu làm theo.
\end{luuy}

\begin{vidu}[minh hoạ một tình huống cụ thể]
Hộp ví dụ tự đánh số theo chương, tiện khi em cần dẫn lại một ví dụ ở chỗ
khác trong báo cáo.
\end{vidu}
